\documentclass[hidelinks,11pt,final]{article}
\usepackage{custom}

\renewcommand{\restriction}{\mathord{\upharpoonright}}
%\newcommand*{\Perm}[2]{{}^{#1}\!P_{#2}}%
%\newcommand*{\Comb}[2]{{}^{#1}C_{#2}}%


\hypersetup{
    pdftitle={Cryptographic perfect hash functions: A theoretical analysis on space efficiency and entropy},                    % title
    pdfauthor={Alexander Towell},               % author
    pdfsubject={computer science},              % subject of the document
    pdfkeywords={
        perfect hash function,
        random oracle,
        space complexity,
        maximum entropy,
        cryptographic hash function},           % keywords
    colorlinks=false,                           % false: boxed links; true: colored links
    citecolor=green,                            % color of links to bibliography
    filecolor=magenta,                          % color of file links
    urlcolor=green                              % color of external links
}

\title
{
    Rate-distorted cryptographic perfect hash functions\\
    \large A theoretical analysis on space efficiency, time complexity, and entropy
}

\author
{
    Alexander Towell\\
    \texttt{atowell@siue.edu}
}
\date{}

%\PassOptionsToPackage{unicode}{hyperref}
%\PassOptionsToPackage{naturalnames}{hyperref}


\setglossarystyle{list}
\loadglsentries{gls-set}
\loadglsentries{gls-hashing}
\loadglsentries{gls-custom}
\makeglossaries

\begin{document}
\maketitle
\begin{abstract}
We analyze a theoretical \emph{perfect hash function} that has three desirable properties: (1) it is a cryptographic hash function; (2) its \emph{in-place} encoding obtains the theoretical lower-bound for the expected space complexity; and (3) its \emph{in-place} encoding is a random bit string with maximum entropy.
\end{abstract}

\textbf{Keywords:} perfect hash function,
        random oracle, space complexity,
        maximum entropy, cryptographic hash function

\tableofcontents

\printglossary[type=symbols,style=list,nogroupskip]
\printglossary[type=\acronymtype]

\subfile{sections/perf_hash_fn}
\subfile{sections/random_oracle_phf}
\subfile{sections/deterministic}
\subfile{sections/perf_hash_2_lvl_fn}
\subfile{sections/rate_distorted}
\subfile{sections/composition}
\subfile{sections/appendix}
\printglossary

%\listoffigures
%\listofalgorithms

\bibliographystyle{plain}
\bibliography{references}
\end{document}

